\exercisesfor{6}

第 6 节的约定除非另有说明均有效。

¶ 1. 设 g 为半单李代数，ρ 为 g 在 M 上的一个有限维表示。

\begin{enumerate}
\def\labelenumi{(\alph{enumi})}
\item
  若 \(\rho\) 单且非零，则 \(\mathrm{H}^p (\mathfrak{g},\mathbf{M}) = \{0\}\) 对一切 \(p\) 成立。（使用第 3 节命题 1 和练习 12 (j)。）
\item
  对一切 \(\rho\)，有 \(\mathbf{H}^1(\mathfrak{g}, \mathbf{M}) = \{0\}\)。（若 \(\rho\) 单且非零，应用 (a)。若 \(\rho\) 为零，利用 \(\mathfrak{g} = \mathcal{D}\mathfrak{g}\) 这一事实。一般情形，对 \(\rho\) 的维数作归纳：若 N 是 M 的异于 \(\{0\}\) 与 M 的子 g-模，利用正合列：
\end{enumerate}

\[
\mathrm{H} ^ {1} (\mathfrak {g}, \mathrm{N}) \rightarrow \mathrm{H} ^ {1} (\mathfrak {g}, \mathrm{M}) \rightarrow \mathrm{H} ^ {1} (\mathfrak {g}, \mathrm{M} / \mathrm{N})
\]

（此正合列已在第 3 节练习 12 (f) 中建立。）由此重新得出第 2 号的注 2。

\begin{enumerate}
\def\labelenumi{(\alph{enumi})}
\setcounter{enumi}{2}
\item
  由 (b) 推出定理 2 的一个证明。（使用第 3 节练习 12 (g)。）再由 (b) 推出 g 的每个导子都是内导子。（使用第 3 节练习 12 (h)。）
\item
  对一切 \(\rho\)，有 \(\mathrm{H}^2(\mathfrak{g}, \mathrm{M}) = \{0\}\)。（仿照 (b) 的论证，只需考虑 \(\rho = 0\) 的情形。设 \(c \in \mathrm{H}^2(\mathfrak{g}, \mathrm{M})\)。按第 3 节练习 12 (i)，考虑中心扩张 \(\mathfrak{h}\)，即 \(\mathfrak{g}\) 被 \(\mathrm{M}\) 的、由 \(c\) 定义的扩张。\(\mathfrak{h}\) 的伴随表示定义了 \(\mathfrak{g}\) 在 \(\mathfrak{h}\) 上的一个表示。由定理 2，\(\mathrm{M}\) 在 \(\mathfrak{h}\) 中有一个在 \(\mathfrak{g}\) 下稳定的补，因而该扩张是平凡的。于是由第 3 节练习 12 (i) 得 \(c = 0\)。）
\item
  由 (b) 与 (d) 推出定理 5 的一个证明。（与正文中一样，可归结为根基交换的情形。）
\end{enumerate}

\begin{enumerate}
\def\labelenumi{\arabic{enumi}.}
\setcounter{enumi}{1}
\item
  设 g 为李代数，r 为其根基，\((a_{0}, a_{1}, \ldots)\) 为 g 的理想序列，定义如下：(1) \(a_{0} = \{0\}\)；(2) \(a_{i+1}/a_{i}\) 是 \(g/a_{i}\) 的一个极大交换理想。设 p 为使 \(a_{p} = a_{p+1} = \cdots\) 的最小整数。证明 \(r = a_{p}\)。（证明 \(g/a_{p}\) 半单。）
\item
  李代数 g 半单的充分必要条件是：在每个包含 g 作为子代数的李代数中 g 都是约化的。（若 g 满足此条件，设 M 为 K 上有限维 g-模。把 M 视为交换代数，作 g 与 M 的半直积，其中 g 是约化的。由此推出 M 半单。）
\item
  设 g 为半单李代数，\(\rho\) 为 g 在有限维空间 M 上的非零单表示。设 h 为相应的半直积。证明 \(h = Dh\)，h 的中心为零，且 h 不是半单代数与可解代数的乘积。
\item
  设 \(\mathfrak{a}\) 为李代数，\(\mathfrak{r}\) 为其根基。若 \(\mathfrak{r}\) 具有特征理想的递降序列 \(\mathfrak{r} = \mathfrak{r}_0 \supset \mathfrak{r}_1 \supset \dots \supset \mathfrak{r}_n = \{0\}\)，使得 \(\dim \mathfrak{r}_i / \mathfrak{r}_{i+1} = 1\) 对 \(0 \leqslant i < n\) 成立，则 \(\mathfrak{g}\) 是半单代数与可解代数的乘积。
\end{enumerate}

（设 s 为 g 的一个 Levi 子代数。对一切 \(x \in s\)，令 \(\rho(x)\) 为 \(ad_{g}x\) 在 r 上的限制。则 \(\rho\) 是一维表示的直和，这些表示都是零表示，因为 \(s = Ds\)。）

\begin{enumerate}
\def\labelenumi{\arabic{enumi}.}
\setcounter{enumi}{5}
\item
  设 \(\mathfrak{g}\) 为李代数，\(\mathcal{D}^{\infty}\mathfrak{g}\) 为诸 \(\mathcal{D}^{p}\mathfrak{g}\)（ \(p = 1,2,\ldots\) ）的交。证明 \(\mathfrak{g}/\mathcal{D}^{\infty}\mathfrak{g}\) 可解，且 \(\mathcal{D}(\mathcal{D}^{\infty}\mathfrak{g}) = \mathcal{D}^{\infty}\mathfrak{g}\)。\(\mathfrak{g}\) 同构于半单代数与可解代数的乘积的充分必要条件是 \(\mathcal{D}^{\infty}\mathfrak{g}\) 半单。
\item
  \begin{enumerate}
  \def\labelenumii{(\alph{enumii})}
  \tightlist
  \item
    设 g 为李代数，b 为 g 的半单理想，a 为 b 在 g 中的中心化子，于是 g 与 \(b \times a\) 恒同。证明对 g 的每个理想 \(\ell\)，有 \(\ell = (\ell \cap \mathfrak{b}) \times (\ell \cap \mathfrak{a})\)。（设 \(\ell_{1}\) 为 \(\ell\) 到 b 上的典范投影；它是 b 的理想，因而 \(\mathcal{D}\ell_{1} = \ell_{1}\)；由此推出 \(\ell_{1} \subset \ell\)。）
  \end{enumerate}
\end{enumerate}

\begin{enumerate}
\def\labelenumi{(\alph{enumi})}
\setcounter{enumi}{1}
\item
  设 \(\beta\) 为 \(\mathfrak{g}\) 上的不变双线性形式。证明 \(\mathfrak{b}\) 与 \(\mathfrak{a}\) 关于 \(\beta\) 正交（利用 \(\mathfrak{b} = \mathcal{D}\mathfrak{b}\) 以及 \(\beta = \beta_{1} + \beta_{2}\) 这一事实，其中 \(\beta_{1}\)（相应地 \(\beta_{2}\) ）是到 \(\mathfrak{a}\)（相应地 \(\mathfrak{b}\) ）上的限制为零的不变双线性形式。
\item
  由 (a) 推出 \(\mathfrak{g}\) 中存在一个最大半单理想。（考虑 \(\mathfrak{g}\) 的一个极大半单理想。）
\end{enumerate}

\begin{enumerate}
\def\labelenumi{\arabic{enumi}.}
\setcounter{enumi}{7}
\tightlist
\item
  设 g 为李代数。若 \(h \neq \{0\}\) 且 g 的包含于 h 中的每个理想都等于 \(\{0\}\) 或 h，则称 g 的理想 h 为极小理想。
\end{enumerate}

\begin{enumerate}
\def\labelenumi{(\alph{enumi})}
\item
  g 的每个单理想都是极小的。
\item
  设 \(\mathfrak{h}\) 为 \(\mathfrak{g}\) 的极小理想，\(\mathfrak{r}\) 为 \(\mathfrak{g}\) 的根基。则或者 \(\mathfrak{h} \subset \mathfrak{r}\)，此时 \(\mathfrak{h}\) 交换；或者 \(\mathfrak{h} \cap \mathfrak{r} = \{0\}\)，此时 \(\mathfrak{h}\) 单。（利用李代数的导理想与半单李代数的单分量都是特征理想这一事实。）
\end{enumerate}

\begin{enumerate}
\def\labelenumi{\arabic{enumi}.}
\setcounter{enumi}{8}
\item
  李代数 g 约化的充分必要条件是其中心 c 等于它的最大幂零理想。（若该条件成立，设 r 为 g 的根基；Dr 包含在 r 的中心中，因而 r 幂零，从而 r = c。）
\item
  设 \(\mathfrak{g}\) 为李代数，使得条件 \(x \in \mathfrak{g}, y \in \mathfrak{g}, [[x, y], y] = 0\) 蕴含 \([x, y] = 0\)。证明 \(\mathfrak{g}\) 约化。（利用第 5 节练习 3 证明根基 \(\mathfrak{r}\)（\(\mathfrak{g}\) 的根基）交换。然后证明 \([\mathfrak{g}, \mathfrak{r}] = 0\)。）
\end{enumerate}

¶ 11. (a) 设 g 为李代数，r 为其根基，s 为 g 的 Levi 子代数，m 为包含 r 的 g 的理想。则存在 s 的理想 t，在 g 中是 m 的补，使得 \([m, t] \subset r\)。

\begin{enumerate}
\def\labelenumi{(\alph{enumi})}
\setcounter{enumi}{1}
\tightlist
\item
  设 g 为李代数，a 为 g 的次不变子代数。则存在合成列 \(g = g_{0} \supset g_{1} \supset \cdots \supset g_{k} = a\)，使得 \(g_{i}\) 是 \(g_{i+1}\) 与一个子代数 \(b_{i}\) 的直和，后者或者是 1 维且包含在 \(r(g_{i})\)（即 \(g_{i}\) 的根基）中，或者是单的，且满足 \([b_{i}, g_{i+1}] \subset r(g_{i+1})\)。（归结为 a 是 g 的理想的情形。设 \(g' = g/a\)，\(r'\) 为 \(g'\) 的根基。代数 \(g'/r'\) 是它的单理想 \(a_{1}, a_{2}, \ldots, a_{c}\) 的乘积。
\end{enumerate}

取 \(\mathfrak{g}'\) 的一个合成列，它由 \(\mathbf{r}'\) 的逐次商均为 1 维的合成列以及理想 \(a_1, a_1 \times a_2, \ldots, a_1 \times a_2 \times \cdots \times a_c\) 的原像组成。然后取该合成列在 \(\mathfrak{g}\) 中的原像。）

\begin{enumerate}
\def\labelenumi{\arabic{enumi}.}
\setcounter{enumi}{11}
\tightlist
\item
  设 g 为李代数，r 为其根基，D 为 g 的一个导子。
\end{enumerate}

\begin{enumerate}
\def\labelenumi{(\alph{enumi})}
\tightlist
\item
  若 \(\mathrm{D}(\mathfrak{r}) = \{0\} \)，则 D 是内导子。（设 c 为 r 在 g 中的中心化子。g 的伴随表示定义了 \(x^{*} \mapsto \rho(x^{*})\)，即 \(g^{*} = g/r\) 在 c 上的一个表示。另一方面，\([\mathrm{D}(\mathfrak{g}), \mathfrak{r}] \subset \mathrm{D}([\mathfrak{g}, \mathfrak{r}]) + [\mathfrak{g}, \mathrm{D}(\mathfrak{r})] = \{0\}\)，因而 \(\mathrm{D}(\mathfrak{g}) \subset c\)，于是 D 定义了一个线性映射 \(D^{*}: g^{*} \to c\)。证明
\end{enumerate}

\[
\mathrm{D} ^ {*} ([ x ^ {*}, y ^ {*} ]) = \rho (x ^ {*}) \mathrm{D} ^ {*} y ^ {*} - \rho (y ^ {*}) \mathrm{D} ^ {*} x ^ {*}
\]

对一切 \(x^{*}, y^{*}\) ∈ \(\mathfrak{g}^{*}\) 成立。由此推出存在 \(a \in \mathfrak{c}\)，使得 \(D^{*}x^{*} = \rho(x^{*})a\) 对一切 \(x^{*} \in \mathfrak{g}^{*}\) 成立（参见~第 2 号，注 2），从而 \(Dx = [x, a]\) 对一切 \(x \in \mathfrak{g}\) 成立。）

\begin{enumerate}
\def\labelenumi{(\alph{enumi})}
\setcounter{enumi}{1}
\tightlist
\item
  若 D 在 r 上与 g 的某个内导子重合，则 D 是内导子。（应用 (a)。）
\end{enumerate}

\begin{enumerate}
\def\labelenumi{\arabic{enumi}.}
\setcounter{enumi}{12}
\item
  设 g 为代数闭域上的李代数，r 为其根基，\(\rho\) 为 g 的有限维单表示。则 \(\rho(z)\) 对 \(z \in r\) 是标量。（归结为 g 约化的情形。此时 r 是 g 的中心。）
\item
  设 X 为 \(\mathfrak{gl}(n, \mathbf{K})\) 中的幂零矩阵。证明在 \(\mathfrak{gl}(n, \mathbf{K})\) 中存在另外两个矩阵 Y、H，使得 \([H, X] = X\)、\([H, Y] = -Y\)、\([X, Y] = H\)。（对 n 作归纳，利用 X 的 Jordan 标准形；于是归结为 \(X = E_{12} + E_{23} + \cdots + E_{n-1,n}\) 的情形；然后取 Y 为诸 \(E_{k,k-1}\) 的线性组合，H 为诸 \(E_{jj'}\) 的线性组合。）
\end{enumerate}

¶ 15. (a) 设 X、Y 为 \(\mathfrak{gl}(n, \mathrm{K})\) 的两个矩阵。证明若 \([X, Y] = X\)，则 X 幂零。（对每个多项式 \(f \in K[T]\)，注意

\[
[ f (X), Y ] = X f ^ {\prime} (X).)
\]

\begin{enumerate}
\def\labelenumi{(\alph{enumi})}
\setcounter{enumi}{1}
\item
  设 \(\mathfrak{g}\) 为李代数，\(h\)、\(x\) 为 \(\mathfrak{g}\) 的两个元素，满足 \([h, x] = x\)，且存在 \(z \in \mathfrak{g}\) 使 \([x, z] = h\)。证明存在 \(y \in \mathfrak{g}\)，使 \([x, y] = h\) 且 \([h, y] = -y\)。（首先注意 \([z, h] + z\) 属于 \(\mathfrak{n}\)，即 \(x\) 在 \(\mathfrak{g}\) 中的中心化子。然后证明 \(\mathfrak{n}\) 在 ad \(h\) 下稳定，且 \(\mathfrak{n}\) 上的 I - ad \(h\) 的限制是双射；为此，注意到 ad \(x\) 幂零（利用 \((a)\)），并证明：若 \(\mathfrak{g}_{\mathfrak{l}}\) 是 \(\mathfrak{g}\) 在 \((\mathrm{ad}x)^{\mathfrak{l}}\) 下的像，则 I - ad \(h\) 在取商后给出 \((\mathfrak{n} \cap \mathfrak{g}_{\mathfrak{l}-1}) / (\mathfrak{n} \cap \mathfrak{g}_{\mathfrak{l}})\) 的一个双射；为此利用关系式 \([\mathrm{ad}z, (\mathrm{ad}x)^{\mathfrak{k}}] = -k(\mathrm{ad}x)^{\mathfrak{k}-1}(\mathrm{ad}h + \frac{k-1}{2}I).)\)）
\item
  设 \(\mathfrak{g}\) 为 \(\mathfrak{gl}(n, \mathbf{K})\) 的子代数；假设对每个幂零矩阵 \(X \in \mathfrak{g}\)，都存在另外两个矩阵 \(H, Y\)（均在 \(\mathfrak{g}\) 中），使得 \([H, X] = X\)、\([H, Y] = -Y\)、\([X, Y] = H\)。设 \(\mathfrak{h}\) 为 \(\mathfrak{g}\) 的子代数，使得存在
\end{enumerate}

h 在 g 中的一个补子空间 m，且满足 \([h, m] \subset m\)。证明 h 对其所有幂零矩阵具有与 g 相同的性质（利用 (b)）。

\begin{enumerate}
\def\labelenumi{\arabic{enumi}.}
\setcounter{enumi}{15}
\item
  设 g 为 \(\mathfrak{gl}(n,\mathbf{K})\) 的子代数，其恒等表示半单。证明每个幂零矩阵 \(X\in g\) 都包含在 g 的一个 3 维单李子代数中。（证明 g 在 \(\mathfrak{gl}(n,\mathbf{K})\) 中约化；然后利用练习 14 与 15 (c)。）
\item
  证明由实单李代数通过将基域从 \(\mathbf{R}\) 扩张到 \(\mathbf{C}\) 而得到的代数并不总是单的。（设 \(\mathfrak{g}\) 为实单李代数。若 \(\mathfrak{g}' = \mathfrak{g}_{(\mathbf{C})}\) 单，设 \(\mathfrak{h}\) 为由 \(\mathfrak{g}'\) 通过把标量域限制到 \(\mathbf{R}\) 而得到的实李代数。我们知道 \(\mathfrak{h}\) 是单的。于是由第 1 节练习 4，\(\mathfrak{h}_{(\mathbf{C})}\) 是两个同构于 \(\mathfrak{g}'\) 的代数的乘积。另参见练习 26 (b)。
\item
  \begin{enumerate}
  \def\labelenumii{(\alph{enumii})}
  \tightlist
  \item
    设 g 为单李代数。g 上的每个不变双线性形式或为零或非退化。若 K 代数闭，则 g 上的每个不变双线性形式 \(\beta\) 都与 Killing 形式 \(\beta_{0}\) 成比例。（考虑向量空间 g 的自同态 \(\sigma\)，它由 \(\beta(x,y)=\beta_{0}(\sigma(x),y)\) 定义；证明 \(\sigma\) 与 \(ad_{g}x\) 交换，从而是标量。）证明当 K 不是代数闭时该结论可能不成立。（利用练习 17 以及不变双线性形式空间的维数在基域扩张下不变这一事实。）
  \end{enumerate}
\end{enumerate}

\begin{enumerate}
\def\labelenumi{(\alph{enumi})}
\setcounter{enumi}{1}
\item
  设 g 为代数闭域上的半单李代数。由 (a) 与练习 7 (b) 推出：g 上不变双线性形式空间的维数等于 g 的单分量个数，且所有这些形式都是对称的。
\item
  设 g 为单李代数，M 为向量空间 g 的对偶空间，带有伴随表示的对偶表示，b 为由 \(x \mapsto x_{M}\) 定义的 g 与 M 的半直积。对 g 中的 y、z 与 M 中的 \(y'\)、\(z'\)，写 \(\beta(y + y', z + z') = \langle y, z' \rangle + \langle z, y' \rangle\)。证明 \(\beta\) 是 b 上的非退化不变对称双线性形式，且不与 b 的任何表示相关联。（注意 b 的根基与幂零根基都等于 M。）
\end{enumerate}

\begin{enumerate}
\def\labelenumi{\arabic{enumi}.}
\setcounter{enumi}{18}
\tightlist
\item
  设 \(g\) 为约化李代数，\(U\) 为其包络代数，\(Z\) 为 \(U\) 的中心，\(V\) 为 \(U\) 中由形如 \(uv - vu\)（ \(u \in U, v \in U\) ）的元素生成的子空间。
\end{enumerate}

\begin{enumerate}
\def\labelenumi{(\alph{enumi})}
\item
  U 是 Z 与 V 的直和。（把第 3 节命题 6 应用于表示 \(x \mapsto \mathrm{ad}_U x\)（\(\mathfrak{g}\) 在 U 上的表示），并注意在第 2 节第 7 号定理 1 的推论 4 的分解 \(\mathrm{U} = \sum_{n} \mathrm{U}^n\) 中，诸 \(\mathrm{U}^n\) 在该表示下稳定。）
\item
  设 \(u \mapsto u^{\natural}\) 为 U 到 Z 上平行于 V 的投影。证明 \((uv)^{\natural} = (vu)^{\natural}, (zu)^{\natural} = zu^{\natural}\) 对一切 \(u \in \mathrm{U}, v \in \mathrm{U}, z \in \mathrm{Z}\) 成立。
\item
  设 R 为 U 的双边理想。证明 \(\mathbf{R} = (\mathbf{R} \cap \mathbf{Z}) + (\mathbf{R} \cap \mathbf{V})\)。（为证 R 的元素 r 在 V 中的分量属于 R，归结为 K 代数闭的情形；注意 R 在 U 上的表示 \(\rho\)（它延拓 \(x \mapsto ad_{U}x\)）下稳定；把 r 分解为它在诸 \(U^{n}\) 中的分量；最后，把《代数》第 VIII 章第 4 节第 2 号定理 1 的推论 2 应用于 \(\rho\) 在 \(U^{n}\) 上的限制。）
\item
  设 \(\mathbf{R}'\) 为 \(Z\) 的理想。设 \(\mathbf{R}_1\) 为 \(U\) 中由 \(\mathbf{R}'\) 生成的双边理想。证明 \(\mathbf{R}_1 \cap Z = \mathbf{R}'\)。（利用 (b)。）
\end{enumerate}

\begin{enumerate}
\def\labelenumi{\arabic{enumi}.}
\setcounter{enumi}{19}
\tightlist
\item
  假设 \(\mathbf{K}\) 代数闭。设 \(\mathfrak{g}\) 为李代数，\(\mathfrak{c}\) 为其中心。
\end{enumerate}

\begin{enumerate}
\def\labelenumi{(\alph{enumi})}
\item
  若 \(\mathfrak{g}\) 容许有限维忠实的单表示，则 \(\mathfrak{g}\) 约化且 \(\dim \mathfrak{c} \leqslant 1\)。
\item
  反之，若 g 约化且 dim \(c \leqslant 1\)，则 g 容许有限维忠实的单表示。（若 g 单，用伴随表示。若 \(g = g_{1} \times g_{2}\) 且 \(g_{1}, g_{2}\) 分别容许有限维忠实的单表示 \(\rho_{1}, \rho_{2}\)（在空间 \(M_{1}, M_{2}\) 上），证明
\end{enumerate}

\[
(x _ {1}, x _ {2}) \mapsto \rho_ {1} (x _ {1}) \otimes 1 + 1 \otimes \rho_ {2} (x _ {2})
\]

是半单表示 \(\rho\)（\(g\) 的），再证明 \(g\) -模 \(\mathrm{M}_1 \otimes \mathrm{M}_2\) 的中心化子归结为标量，从而 \(\rho\) 单。当 \(g_1\) 与 \(g_2\) 都半单，或 \(g_1\) 半单且 \(g_2\) 交换时，通过考察 \(\rho\) 的核证明 \(\rho\) 忠实（该核是 \(g_1 \times g_2\) 的理想）。

\begin{enumerate}
\def\labelenumi{\arabic{enumi}.}
\setcounter{enumi}{20}
\item
  设 V 为 K 上有限维向量空间，维数为 n \textgreater{} 1。证明 \(\mathfrak{sl}(V)\) 单。（归结为 K 代数闭的情形。设 \(\mathfrak{sl}(V) = a \times b\)，其中 dim a = a \textgreater{} 0，dim b = b \textgreater{} 0。设 A（相应地 B）为由 1 与 a（相应地 b）生成的结合代数。于是 V 可视为单 \((\mathrm{A} \otimes_{\mathrm{K}} \mathrm{B})\) -模。因而存在 K 上有限维 p 的 A-模 P 与 K 上有限维 q 的 B-模 Q，使得 V \((\mathrm{A} \otimes_{\mathrm{K}} \mathrm{B})\) -同构于 \(P \otimes_{K} Q\)（《代数》第 VIII 章第 7 节第 7 号命题 8 及第 4 号定理 2）；P 与 Q 忠实，因而 \(a \leqslant p^{2}\)、\(b \leqslant q^{2}\)。另一方面，\(a + b = n^{2} - 1\) 且 pq = n，于是 \((p^{2} - 1)(q^{2} - 1) \leqslant 2\)，矛盾。）
\item
  \begin{enumerate}
  \def\labelenumii{(\alph{enumii})}
  \tightlist
  \item
    设 \(\mathfrak{g}\) 为任意特征的代数闭域 \(\mathbf{K}\) 上的幂零李代数。设 \(\mathfrak{z}\) 为其中心，\(\rho\) 为 \(\mathfrak{g}\) 的有限维表示，\(\beta\) 为相应的双线性形式。证明 \(\mathfrak{z}\cap \mathcal{D}\mathfrak{g}\) 与 \(\mathfrak{g}\) 关于 \(\beta\) 正交。（利用 Jordan-Hölder 列归结为 \(\rho\) 单的情形。然后注意，对 \(z\in \mathfrak{g}\cap \mathcal{D}\mathfrak{g},\rho (z)\) 是迹为零的标量矩阵。再利用第 4 节练习 22 (b)。）由此推出，若 \(\beta\) 非退化，则 \(\mathfrak{g}\) 交换。（利用第 4 节练习 3 (a)。）
  \end{enumerate}
\end{enumerate}

\begin{enumerate}
\def\labelenumi{(\alph{enumi})}
\setcounter{enumi}{1}
\item
  假设 \(\mathbf{K}\) 的特征为 2。证明对可解李代数 \(\mathfrak{gl}(2,\mathbf{K}) = \mathfrak{g}\)，与恒等表示相应的双线性形式非退化，但中心 \(\mathfrak{z}\) 包含在 \(\mathcal{D}\mathfrak{g}\) 中且 \(\neq \{0\}\)。
\item
  设 g 为 K 上以 \((a, b, c, d, e, f)\) 为基的 6 维李代数，其乘法表为 \([a, b] = -[b, a] = d\)、\([a, c] = -[c, a] = e\)、\([b, c] = -[c, b] = f\)，其余括积均为 0。设 β 为 g 上的双线性形式，满足 β \((c, d) = \beta(d, c) = 1\)、β \((a, f) = \beta(f, a) = 1\)、β \((b, e) = \beta(e, b) = -1\)，且 β 在所论基的有序偶上的其余取值均为 0。证明 β 不变，g 幂零，\(\mathfrak{z} = \mathcal{D}g \neq \{0\}\)，且 β 非退化。
\end{enumerate}

\begin{enumerate}
\def\labelenumi{\arabic{enumi}.}
\setcounter{enumi}{22}
\tightlist
\item
  设 \(\mathfrak{g}\) 为任意特征的域 \(\mathbf{K}\) 上的 3 维非交换李代数。
\end{enumerate}

\begin{enumerate}
\def\labelenumi{(\alph{enumi})}
\item
  若 \(\mathfrak{g}\) 的中心 \(\mathfrak{z} \neq \{0\}\)，则 dim \(\mathfrak{z} = 1\)（第 1 节练习 2）且 dim \(\mathcal{D}_{\mathfrak{g}} = 1\)。若 \(\mathfrak{z} \neq \mathcal{D}_{\mathfrak{g}}, \mathfrak{g}\) 是 \(\mathfrak{z}\) 与 2 维非交换代数的乘积。若 \(\mathfrak{z} = \mathcal{D}_{\mathfrak{g}}, \mathfrak{g}\) 是 3 维非交换幂零代数（第 4 节练习 9 (a)）。
\item
  若 \(\mathfrak{z} = \{0\}\) 但 \(\mathfrak{g}\) 中存在 2 维交换子代数，则该子代数唯一且等于 \(\mathcal{D}_{\mathfrak{g}}\)；对一切 \(x \notin \mathcal{D}_{\mathfrak{g}}\)，\(u\)（\(ad x\) 在 \(\mathcal{D}_{\mathfrak{g}}\) 上的限制）是该向量空间的一个双射，且在相差一个乘法常数下唯一确定。反之，每个自同构 \(u\)（2 维向量空间 \(\mathrm{Ka} + \mathrm{Kb}\) 的）都在 \(\mathfrak{g} = \mathrm{Ka} + \mathrm{Kb} + \mathrm{Kc}\) 上由条件 \([a, b] = 0\)、\([c, a] = u(a)\)、\([c, b] = u(b)\) 确定一个李代数结构。两个如此由自同构 \(u_1, u_2\) 定义的李代数同构的充分必要条件是 \(u_1\) 与 \(u_2\) 的矩阵在相差一个标量因子下相似。
\item
  假设 g 中不存在维数 \textgreater1 的交换子代数。证明此时存在 \(a \in g\)，使 \(a \notin (\operatorname{ad} a)g\)（假设某元素 x 不具有此性质，利用 Jacobi 恒等式证明另一元素具有所需性质）。于是存在 g 的基 \((a, b, c)\)，使 \([a, b] = c\)、\([a, c] = \beta b\)、\([b, c] = \gamma a\)，其中 \(\beta \neq 0\) 且 \(\gamma \neq 0\)，且 g 单。设 \(\phi\) 为 g 到外幂 \(\bigwedge^{2}\mathfrak{g}^{*}\)（\(\mathfrak{g}\) 的对偶空间的外幂）上的典范同构，它在相差一个可逆常数因子下确定（《代数》第 III 章第 8 节第 5 号）；设 \(u\) 为 \(\mathfrak{g}\) 到 \(\mathfrak{g}^{*}\) 的、由条件 \(\langle [x,y],u(z)\rangle = \langle x\wedge y,\phi (z)\rangle\) 定义的线性映射；双线性形式 \(\Phi (x,y) = \langle x,u(y)\rangle\)（在 \(\mathfrak{g}\times \mathfrak{g}\) 上）对称且非退化，且 \(2\Phi\) 是 \(\mathfrak{g}\) 的 Killing 形式，相差一个 \(\neq 0\) 的因子。K 上两个 3 维单李代数同构的充分必要条件是相应的双线性形式在相差一个 \(\neq 0\) 的常数因子下等价。
\item
  若 \(\mathbf{K}\) 的特征不是 2，则 (c) 中定义的单代数 \(\mathfrak{g}\) 同构于李代数 \(\mathfrak{o}(\Phi)\)（即形式正交群 \(\mathbf{G}(\Phi)\) 的李代数）（第 1 节练习 26 (a)）。要在 \(\mathfrak{g}\) 中存在向量 \(x\)，使 ad \(x\) 容许与 \(x\) 不成比例的特征向量，其充分必要条件是 \(\Phi\) 的指标 \(>0\)；此时 \(\mathfrak{g}\) 容许一个基 \((a, b, c)\)，使 \([a, b] = b\)、\([a, c] = -c\)、\([b, c] = a\)。
\item
  若 K 的特征为 2，证明 g 上不存在 2-映射，因而 g 不是某个形式群的李代数。证明 g 容许非内导子的导子。
\end{enumerate}

¶ 24. 设 K 为任意特征 p 的域。~我们再次采用定义 1。

\begin{enumerate}
\def\labelenumi{(\alph{enumi})}
\tightlist
\item
  证明除非 \(p = n = 2\)，\(\mathfrak{gl}(n, \mathbf{K})\) 的仅有的非平凡理想
\end{enumerate}

是 1 维中心 \(\mathfrak{z}\)（以 \(\sum_{i=1}^{n} E_{ii}\) 为基）以及由迹为 0 的矩阵组成的李代数 \(\mathfrak{sl}(n, \mathrm{K})\)。（除所指出的例外情形外，注意若理想 \(\mathfrak{a}\) 包含某个 \(E_{ij}\)，则它必然包含 \(\mathfrak{sl}(n, \mathrm{K})\)。若 \(\mathfrak{a}\) 包含一个不属于 \(\mathfrak{b}\) 的元素，则用至多四个适当选取的元素 \(E_{ij}\) 去乘该元素，即得到某个 \(E_{ij}\) 的非零倍数。）若 \(n\) 不是 \(p\) 的倍数，证明 \(\mathfrak{gl}(n, \mathrm{K})\) 是 \(\mathfrak{z}\) 与 \(\mathfrak{sl}(n, \mathrm{K})\) 的直和，且 \(\mathfrak{sl}(n, \mathrm{K})\) 单。反之，若 \(n\) 是 \(p\) 的倍数且 \(n > 2\)，证明 \(\mathfrak{sl}(n, \mathrm{K}) / \mathfrak{z}\) 单（同样的方法）。

\begin{enumerate}
\def\labelenumi{(\alph{enumi})}
\setcounter{enumi}{1}
\tightlist
\item
  证明当 \(n\) 是 \(p\) 的倍数且 \(n > 2\) 时，\(\mathfrak{gl}(n, \mathbf{K}) / \mathfrak{sl}(n, \mathbf{K})\) 的根基为 \(\{0\}\)，但它有交换的商 \(\mathfrak{gl}(n, \mathbf{K}) / \mathfrak{sl}(n, \mathbf{K})\)。
\end{enumerate}

¶ 25. 设 K 为任意特征 p 的域。~设 \(\mathfrak{sp}(2n, \mathrm{K})\) 为 K 上 2n 个变量的形式辛群的李代数（第 1 节练习 26 (c)）。证明该李代数（它与 \(\mathfrak{gl}(2n, \mathrm{K})\) 的一个子代数恒同）具有由下列元素组成的基

\[
\begin{array}{c} H _ {i} = E _ {i i} - E _ {i + n, i + n} \quad (1 \leqslant i \leqslant n), \\ F _ {i j} = E _ {i j} - E _ {j + n, i + n} \quad (1 \leqslant i, j \leqslant n, i \neq j), \\ G _ {i j} ^ {\prime} = E _ {i, j + n} + E _ {j, i + n}, \qquad G _ {i j} ^ {\prime \prime} = E _ {i + n, j} + E _ {j + n, i} \quad (1 \leqslant i <   j \leqslant n), \\ E _ {i, i + n} \quad \text { and } \quad E _ {i + n, i} \quad (1 \leqslant i \leqslant n). \end{array}
\]

\begin{enumerate}
\def\labelenumi{(\alph{enumi})}
\item
  证明若 \(p \neq 2\) 且 \(n \geqslant 1\)，则代数 \(\mathfrak{sp}(2n, \mathrm{K})\) 单（当 \(n = 1\) 时，\(\mathfrak{sp}(2, \mathrm{K}) = \mathfrak{sl}(2, \mathrm{K})\)）。（与练习 24 同法。）
\item
  若 \(p = 2\) 且 \(n \geqslant 3\)，证明 \(H_{i}, F_{ij}, G_{ij}'\) 与 \(G_{ij}''\) 构成理想 \(\mathfrak{a}\)，维数为 \(n(2n - 1)\)；\(\mathfrak{a}\) 中形如
\end{enumerate}

\[
\sum_ {i = 1} ^ {n} \lambda_ {i} H _ {i} + \sum_ {i \neq j} \alpha_ {i j} F _ {i j} + \sum_ {i <   j} \beta_ {i j} G _ {i j} ^ {\prime} + \sum_ {i <   j} \gamma_ {i j} G _ {i j} ^ {\prime \prime}
\]

且满足 \(\sum_{i=1}^{n} \lambda_i = 0\) 的元素构成理想 \(b \subset a\)，维数为 \(n(2n - 1) - 1\)，而 \(\sum_{i=1}^{n} H_i\) 的倍数构成一个理想 \(c\)，即 \(\mathfrak{sp}(2n, K)\) 的中心；\(b, c\) 是 \(\mathfrak{a},\mathfrak{b} = \mathcal{D}(\mathfrak{sp}(2n,\mathbf{K}))\) 仅有的理想，且 \(\mathfrak{b} / \mathfrak{c}\) 单。（同法。）\(\mathfrak{sp}(4,\mathbf{K})\) 在 \(\mathbf{K}\) 的特征为 2 时的理想是什么？

¶ 26. 设 K 为特征 ≠2 的域，Φ 为 K 上 n 维向量空间上的非退化对称双线性形式。

\begin{enumerate}
\def\labelenumi{(\alph{enumi})}
\tightlist
\item
  证明若 \(n \geqslant 5\)，则李代数 \(\mathfrak{o}(\Phi)\) 单。（通过扩张基域，可归结为 \(\Phi\) 的指标为 \(m = [n/2]\) 的情形。例如当 \(n = 2m\) 为偶数时，\(\mathfrak{o}(\Phi)\) 具有由下列元素组成的基
\end{enumerate}

\[
\begin{array}{c} H _ {i} = E _ {i i} - E _ {i + m, i + m}, \quad F _ {i j} = E _ {i j} - E _ {j + m, i + m} \quad (1 \leqslant i, j \leqslant m, i \neq j), \\ G _ {i j} ^ {\prime} = E _ {i, j + m} - E _ {j, i + m} \quad \text { and } \quad G _ {i j} ^ {\prime \prime} = E _ {i + m, j} - E _ {j + m, i} \quad (1 \leqslant i <   j \leqslant m); \end{array}
\]

然后仿照练习 24 与 25 论证。当 \(n = 2m + 1\) 为奇数时同样处理。）

\begin{enumerate}
\def\labelenumi{(\alph{enumi})}
\setcounter{enumi}{1}
\tightlist
\item
  设 \(\Delta\) 为 \(\Phi\) 关于某个基的判别式。设 \(n = 4\)。证明若 \(\Delta\) 在 \(\mathbf{K}\) 中不是平方元，则 \(\mathfrak{o}(\Phi)\) 单。若 \(\Delta\) 是平方元，则 \(\mathfrak{o}(\Phi)\) 是两个互相同构的 3 维单李代数的乘积。（利用《代数》第 IX 章第 9 节练习 16 中描述的 \(\mathbf{G}(\Phi)\) 的结构。）由此给出一个单李代数在扩张基域后变为非单的例子。
\end{enumerate}

\begin{enumerate}
\def\labelenumi{\arabic{enumi}.}
\setcounter{enumi}{26}
\tightlist
\item
  \begin{enumerate}
  \def\labelenumii{(\alph{enumii})}
  \tightlist
  \item
    证明在有限维李 \(p\) -代数中，根基是 \(p\) -理想。（利用第 1 节练习 22 (c)。）
  \end{enumerate}
\end{enumerate}

\begin{enumerate}
\def\labelenumi{(\alph{enumi})}
\setcounter{enumi}{1}
\tightlist
\item
  李 \(p\) -代数 \(\mathfrak{g}\) 的根基为零的充分必要条件是 \(\mathfrak{g}\) 不含交换 \(p\) -理想 \(\neq \{0\}\)。（利用第 1 节练习 22 (c)。）
\end{enumerate}

